% ============================================================
%  MEMORIA FORMAL — Optimización DxI-9000
%  Compilar con: pdflatex o xelatex (Overleaf: pdflatex)
% ============================================================
\documentclass[12pt, a4paper]{article}

% ─── Paquetes ──────────────────────────────────────────────
\usepackage[T1]{fontenc}
\usepackage[utf8]{inputenc}
\usepackage[spanish,es-noshorthands]{babel}
\usepackage{lmodern}
\usepackage{microtype}

\usepackage{geometry}
\geometry{left=2.5cm, right=2.5cm, top=3cm, bottom=3cm}

\usepackage{amsmath, amssymb, amsthm}
\usepackage{mathtools}
\usepackage{booktabs}
\usepackage{array}
\usepackage{longtable}
\usepackage{multirow}
\usepackage{tabularx}
\usepackage{xcolor}
\usepackage{graphicx}
\usepackage{float}
\usepackage{caption}
\usepackage{subcaption}
\usepackage{enumitem}
\usepackage{hyperref}
\usepackage{listings}
\usepackage{parskip}
\usepackage{fancyhdr}
\usepackage{titlesec}
\usepackage{pifont}

% ─── Hiperenlaces ──────────────────────────────────────────
\hypersetup{
  colorlinks=true,
  linkcolor=black,          % índice y referencias internas en negro
  citecolor=black,
  urlcolor=blue!60!black,   % URLs externas conservan el azul
  pdftitle={Optimización de la distribución de reactivos en analizadores DxI-9000},
  pdfauthor={Laboratorio de Optimización},
}

% ─── Cabecera y pie de página ──────────────────────────────
\pagestyle{fancy}
\fancyhf{}
\fancyhead[L]{\small Laboratorio de Optimización}
\fancyhead[R]{\small DxI-9000 — Distribución de reactivos}
\fancyfoot[C]{\thepage}
\renewcommand{\headrulewidth}{0.4pt}

% ─── Código fuente ─────────────────────────────────────────
\lstset{
  basicstyle=\ttfamily\small,
  breaklines=true,
  frame=single,
  backgroundcolor=\color{gray!8},
  rulecolor=\color{gray!40},
  xleftmargin=1em,
  xrightmargin=1em,
}

% ─── Checks y cruces ───────────────────────────────────────
\newcommand{\cmark}{\ding{51}}
\newcommand{\xmark}{\ding{55}}

% ============================================================
\begin{document}
% ============================================================

% ─── PORTADA ───────────────────────────────────────────────
\begin{titlepage}
\thispagestyle{empty}
\centering

% Espacio para el logo de la Universidad de Murcia
% Descomenta la línea siguiente y ajusta el nombre del archivo
% una vez que hayas subido el logo a Overleaf:
% \includegraphics[width=5cm]{logo_umu.png}\\[1cm]
\vspace*{1.5cm}
\rule{\linewidth}{0.5pt}\\[0.4cm]

{\fontsize{22}{26}\selectfont\bfseries
  Optimización de la distribución\\[6pt]
  de reactivos en analizadores DxI-9000}\\[0.3cm]

\rule{\linewidth}{0.5pt}\\[1.8cm]

{\large\itshape Trabajo de la asignatura}\\[0.4cm]
{\LARGE\bfseries Laboratorio de Optimización}\\[0.3cm]
{\large Curso 2025/2026}\\[2cm]

{\large\bfseries Universidad de Murcia}\\[0.3cm]
{\normalsize Facultad de Informática}\\[2.5cm]

\begin{tabular}{rl}
  \textbf{Autores:} & Álvaro Espejo Martínez \\
                    & Javier Hernández Rosique \\
                    & Raúl Sánchez Ibáñez \\[1.2em]
  \textbf{Datos:}   & Enero 2026 --- Laboratorio de Análisis Clínico \\
                    & Hospital General Universitario Morales Meseguer \\
\end{tabular}

\vfill
{\small Murcia, \the\year}

\end{titlepage}

% ─── Agradecimientos ───────────────────────────────────────
\clearpage
\thispagestyle{empty}
\vspace*{3cm}
{\Large\bfseries Agradecimientos}
\vspace{1.5em}

\noindent
Queremos expresar nuestro más sincero agradecimiento a todos los miembros del
Laboratorio de Análisis Clínicos del Hospital General Universitario Morales Meseguer
y, en especial, a \textbf{Francisco Espejo López}, por abrirnos las puertas de sus
instalaciones, recibirnos con tal hospitalidad y mostrarse siempre disponible para
colaborar y ser partícipe de la elaboración del presente proyecto.
Su generosidad y dedicación han sido fundamentales para que este trabajo pudiera
llevarse a cabo con rigor y con una comprensión real del entorno clínico que lo motiva.

\vfill

% ─── Resumen (Abstract) ────────────────────────────────────
\clearpage
\thispagestyle{empty}
\vspace*{1.5cm}
{\Large\bfseries Resumen}
\vspace{1em}

\noindent
Este trabajo presenta un modelo de programación lineal entera mixta (MILP) para optimizar
la distribución de reactivos en los tres analizadores DxI-9000 de un laboratorio de
análisis clínico. El problema consiste en decidir en qué analizadores instalar cada una
de las 28 pruebas clínicas y cuántas petacas de reactivo asignar a cada posición, con el
doble objetivo --- jerárquico y lexicográfico --- de minimizar el tiempo máximo de
procesamiento diario (\emph{makespan}) y maximizar la co-ubicación de grupos de pruebas
que se solicitan conjuntamente en la misma petición clínica.

\medskip\noindent
El modelo se formula en dos fases secuenciales. La Fase 1 minimiza el makespan, logrando
un balance perfecto entre los tres analizadores (172.2\,min, 35.9\,\% de la jornada).
La Fase 2, con el makespan fijado en su valor óptimo, maximiza la cohesión jerárquica de
peticiones mediante variables binarias que capturan la co-ubicación de grupos de hasta
seis pruebas, linealizadas con la generalización de McCormick para $n$ binarias. El modelo
cuenta con 403 variables y 1\,188 restricciones y es resuelto con el solver HiGHS 1.11.0.

\medskip\noindent
La solución óptima reduce el makespan un 8.7\,\% respecto a la distribución empírica
actual, elimina el desequilibrio de carga entre equipos y alcanza una cohesión del
99.7\,\% (25\,804 de 25\,872 co-ocurrencias mensuales cubiertas en un único analizador).
Un postproceso proporcional garantiza el llenado completo de las 144 posiciones de
reactivo con factores de cobertura mínimos de 2.0$\times$ la demanda máxima en todas las
pruebas.

\medskip
\noindent\textbf{Palabras clave:} MILP, optimización de laboratorio, analizadores clínicos,
makespan, cohesión de peticiones, linealización AND, DxI-9000.

\vfill

% ─── Índice ────────────────────────────────────────────────
\clearpage
\thispagestyle{empty}
\tableofcontents

\clearpage
% ============================================================
\section{Contexto del laboratorio}
% ============================================================

\subsection{¿Qué hace este laboratorio?}

El \textbf{Laboratorio de Análisis Clínico} recibe cada día cientos de muestras de sangre
procedentes de pacientes hospitalizados y de consultas externas. Su misión es medir la
concentración de hormonas, marcadores tumorales y vitaminas en el suero de esas muestras,
y devolver los resultados al médico solicitante en el menor tiempo posible.

El proceso comienza cuando un médico emite una \textbf{petición}: un documento —hoy
completamente digitalizado— que especifica qué compuestos hay que medir en la muestra de
un paciente concreto. Una petición puede contener desde una sola determinación
(p.\,ej., ``medir TSH'') hasta diecisiete pruebas distintas; la mediana observada en enero
de 2026 es de \textbf{2 determinaciones por petición}.

Las muestras se obtienen centrifugando sangre venosa extraída en tubos de tapón rojo con
gel separador. El centrifugado separa el suero del coágulo sólido; el suero es lo que
realmente analizan los equipos. Cada tubo queda asociado a una petición concreta, con una
lista de pruebas pendientes que deben realizarse antes de emitir el informe clínico.

\subsection{Del paciente al resultado: el recorrido de un tubo}

Una vez centrifugado, el tubo entra en la \textbf{cadena de distribución automatizada} del
laboratorio: una cinta transportadora que lo identifica por código de barras, registra sus
pruebas pendientes y lo dirige hacia los analizadores correctos sin intervención humana.

Los tres analizadores \textbf{DxI-9000} están dispuestos en serie a lo largo de la cadena,
formando una línea de procesamiento continua:

\begin{center}
\ttfamily\small
\begin{tabular}{c}
Recepción/centrifugado \\
$\downarrow$ \\
{[CER --- identificación]} \\
$\downarrow$ \\
{[DxI-9000 n.º 1] $\longrightarrow$ [DxI-9000 n.º 2] $\longrightarrow$ [DxI-9000 n.º 3]} \\
$\downarrow$ \\
Resultados $\rightarrow$ Historia clínica
\end{tabular}
\end{center}

Cuando un tubo llega frente a un DxI-9000, el analizador consulta la lista de pruebas
pendientes: si \textbf{tiene instalada} alguna de esas pruebas, retiene el tubo y la
procesa; en caso contrario, el tubo pasa de largo. Un tubo puede ser retenido por uno,
dos o los tres analizadores dependiendo de la distribución de sus pruebas.
\textbf{Este reparto de pruebas entre analizadores es exactamente la decisión que el modelo
de optimización debe tomar.}

\subsection{Los tres DxI-9000: iguales pero independientes}

Los tres equipos son del mismo modelo (Beckman Coulter DxI-9000) y realizan técnicamente
las mismas pruebas. Sin embargo, en cada momento pueden tener instaladas pruebas distintas
en cantidades distintas y procesar colas de tubos independientes.

Una diferencia operativa relevante es la \textbf{capacidad de posiciones de reactivo}:
DxI1 dispone de 50 posiciones, mientras que DxI2 y DxI3 tienen 47 cada uno (las 3
posiciones restantes están permanentemente reservadas para muestras urgentes y no se
incluyen en este modelo).

\clearpage
% ============================================================
\section{Flujo de trabajo y elementos clave}
% ============================================================

\subsection{Glosario esencial}

\begin{table}[H]
\centering
\caption{Terminología clave del laboratorio.}
\begin{tabularx}{\linewidth}{lX}
\toprule
\textbf{Término} & \textbf{Definición} \\
\midrule
Determinación & Medición individual de un compuesto en el suero de un paciente. Unidad mínima de trabajo. \\
Prueba & Tipo de análisis (p.\,ej., TSH, VITD, PSA). Cada prueba requiere su propio reactivo. \\
Petición & Solicitud médica que agrupa varias determinaciones para un mismo paciente. \\
Tubo & Recipiente físico con la muestra de suero. Lleva asociada una petición. \\
Reactivo & Sustancia química específica que produce la reacción medible al mezclarse con el suero. \\
Petaca de reactivo & Cartucho sellado con reactivo para un número limitado de determinaciones (50--200). \\
Posición de reactivo & Hueco en el carrusel del analizador donde se aloja una petaca. \\
Paso (\textit{step}) & Fase de incubación del ensayo. 1 paso: 450\,det/h; 2 pasos: 225\,det/h. \\
Makespan & Tiempo del analizador más cargado en procesar toda su cola del día. \\
Carro de cubetas & Componente interno con 269 posiciones dinámicas para las alícuotas de suero. \\
\bottomrule
\end{tabularx}
\end{table}

\subsection{Cómo procesa un tubo el DxI-9000}

Cada DxI-9000 es un \textbf{analizador de inmunoensayo quimioluminiscente}. El ciclo de
una determinación incluye: (1) identificación del tubo, (2) aspiración de la alícuota,
(3) incubación con el reactivo (1 o 2 fases), (4) medición de la señal, y (5) volcado
automático al sistema de información (LIS).

\begin{itemize}
  \item \textbf{Ensayos de 1 paso:} hasta $\mathbf{450}$\,\textbf{determinaciones/hora}.
  \item \textbf{Ensayos de 2 pasos:} hasta $\mathbf{225}$\,\textbf{determinaciones/hora}.
\end{itemize}

El cuello de botella real no es la velocidad del analizador sino la
\textbf{disponibilidad de reactivo}: si una petaca se agota a mitad de jornada, el equipo
queda bloqueado esperando recarga, que además requiere calibración y control de calidad.
Por ello el número de petacas es la variable de decisión más crítica.

\subsection{El recurso escaso: las posiciones de reactivo}

\begin{itemize}
  \item \textbf{DxI1:} 50 posiciones disponibles.
  \item \textbf{DxI2 y DxI3:} 47 posiciones disponibles cada uno.
  \item \textbf{Total del sistema:} 144 posiciones para repartir entre 28 pruebas y 3 analizadores.
\end{itemize}

El DxI-9000 selecciona la petaca por cercanía física en el carrusel; tener posiciones
vacías penaliza el tiempo de ciclo. Por ello el laboratorio exige que
\textbf{todas las posiciones estén siempre ocupadas}.

\subsection{Las 28 pruebas clínicas}

\begin{table}[H]
\centering
\caption{Clasificación técnica de las 28 pruebas.}
\begin{tabular}{llll}
\toprule
\textbf{Tipo} & \textbf{Velocidad} & \textbf{Det/petaca} & \textbf{Pruebas} \\
\midrule
1 paso & 450\,det/h & 50--200$^*$ & TSH, VITD, PSA, PSAL, CA125, AFP, PTH, \\
       &            &             & T3L, PRL, HCG, TES, INS, PEPC, COR, TIRO \\
2 pasos & 225\,det/h & 50 & B12, FOL, T4L, CEA, CA199, CA153, \\
        &            &    & FSH, LH, EST, PROG, TPO, ATG, SHBG \\
\bottomrule
\multicolumn{4}{l}{$^*$TSH: 200\,det/petaca; PSA: 100\,det/petaca; resto: 50\,det/petaca.}
\end{tabular}
\end{table}

Las pruebas se agrupan clínicamente en tres paneles:
\begin{itemize}
  \item \textbf{Grupo A — Tiroideo-metabólico (alta demanda):} TSH, T4L, T3L, TPO, VITD, B12, FOL.
  \item \textbf{Grupo B — Marcadores tumorales + PSA:} CEA, CA125, CA153, CA199, PSA, PSAL, AFP.
  \item \textbf{Grupo C — Panel hormonal reproductor:} FSH, LH, EST, PRL, PROG, TES, HCG, INS, COR, PTH, ATG, TIRO, PEPC, SHBG.
\end{itemize}

\clearpage
% ============================================================
\section{Descripción del problema de optimización}
% ============================================================

\subsection{Objetivo}

El modelo persigue dos objetivos de forma \textbf{jerárquica (lexicográfica)}:

\begin{enumerate}
  \item \textbf{Fase 1 — Minimizar el makespan:} distribuir las 28 pruebas entre los 3
        analizadores y decidir cuántas petacas asignar a cada una, de forma que se minimice
        el tiempo máximo de procesamiento diario equilibrando la carga.
  \item \textbf{Fase 2 — Maximizar la cohesión de peticiones:} manteniendo el makespan
        óptimo, redistribuir la asignación para que los grupos de pruebas que aparecen
        juntos con mayor frecuencia queden instalados en el mismo analizador.
\end{enumerate}

\subsection{Decisiones a tomar}

\begin{enumerate}
  \item \textbf{¿En qué analizadores se instala cada prueba?} Una prueba puede estar en uno,
        dos o tres analizadores simultáneamente.
  \item \textbf{¿Cuántas petacas de cada prueba se cargan en cada analizador?}
  \item \textbf{¿Qué grupos de pruebas quedan co-ubicados?} Las variables $w_g^{(\ell)}$
        capturan si el grupo $g$ de nivel $\ell$ está íntegramente en el analizador $j$.
\end{enumerate}

\subsection{Factores clave}

\paragraph{Cuello de botella por exceso de carga.}
Si un analizador acumula demasiadas pruebas de alta demanda, su cola crece y el tiempo
total de procesamiento se dispara mientras los otros equipos están infrautilizados.

\paragraph{Cuello de botella por escasez de petacas.}
Si una prueba de alta demanda tiene pocas petacas, el reactivo se agota frecuentemente y
el analizador debe esperar a que se reponga, bloqueando el procesamiento. Cada recarga
exige además calibración y control de calidad, consumiendo decenas de minutos adicionales.

\paragraph{Fragmentación de peticiones.}
Si las pruebas de una misma petición están repartidas entre dos analizadores, el tubo de
suero tiene que pasar por ambos, consumiendo el doble de recursos del carro de cubetas.
Mantener grupos de pruebas que se solicitan juntas en el mismo analizador reduce esta
fragmentación y acelera la entrega de resultados al médico.

\subsection{Restricciones del problema}

\begin{table}[H]
\centering
\caption{Restricciones del modelo.}
\begin{tabularx}{\linewidth}{clX}
\toprule
\textbf{\#} & \textbf{Restricción} & \textbf{Motivo} \\
\midrule
R1  & Petacas por analizador $\leq$ posiciones disponibles & Límite físico del carrusel \\
R2  & Petacas suficientes para la demanda enrutada & Evitar agotamiento de reactivo \\
R3  & Toda la demanda diaria cubierta & Ninguna determinación sin procesar \\
R4  & Enrutar solo si la prueba está instalada & Restricción lógica \\
R5  & Coherencia petacas $\leftrightarrow$ presencia & Consistencia del modelo \\
R6  & Mayoría de pruebas en $\geq 2$ analizadores & Contingencia operativa \\
R6b & Big Six en los \textbf{3 analizadores} & Pruebas de mayor volumen \\
R6c & Pruebas con dem. $< 20$\,det/día en $\leq 2$ analizadores & Coste calibración $>$ beneficio \\
R7  & PSAL $\leq$ PSA por analizador & PSAL siempre con PSA \\
R8  & Makespan $\geq$ carga de cada analizador & Definición del makespan \\
R9  & $|\text{pruebas 1 paso}_j - \text{pruebas 1 paso}_k| \leq \delta$ & Equilibrio de velocidades \\
R10 & Makespan $\leq 8$\,h & Jornada laboral \\
R11 & Makespan $\leq T^*$ (solo Fase 2) & No empeorar la Fase 1 \\
R15--R16 & Linealización de $w_g^{(\ell)}$ (AND de $n$ binarias) & Reemplaza no-linealidad \\
\bottomrule
\end{tabularx}
\end{table}

\clearpage
% ============================================================
\section{Análisis exploratorio de los datos}
% ============================================================

\subsection{Demanda diaria por prueba (enero 2026, 22 días laborables)}

\begin{table}[H]
\centering
\caption{Estadísticos de demanda por prueba (enero 2026).}
\small
\begin{tabular}{lrrrrrrc}
\toprule
\textbf{Prueba} & \textbf{Pasos} & \textbf{Det/pet.} & \textbf{Media/día} & \textbf{Mín} & \textbf{Máx} & \textbf{CV (\%)} & \textbf{Pet. mín$^*$} \\
\midrule
TSH   & 1 & 200 & 405.4 & 0 & 765 & 78.2  & 3 \\
VITD  & 1 &  50 & 170.8 & 0 & 322 & 77.8  & 4 \\
B12   & 2 &  50 & 153.5 & 0 & 300 & 77.7  & 4 \\
FOL   & 2 &  50 & 126.1 & 0 & 260 & 77.9  & 3 \\
T4L   & 2 &  50 & 118.2 & 0 & 240 & 78.3  & 3 \\
PSA   & 1 & 100 &  77.7 & 0 & 149 & 78.5  & 1 \\
CEA   & 2 &  50 &  41.7 & 0 &  88 & 78.6  & 1 \\
FSH   & 2 &  50 &  21.5 & 0 &  45 & 80.4  & 1 \\
TES   & 1 &  50 &  21.4 & 0 &  43 & 79.5  & 1 \\
CA153 & 2 &  50 &  19.7 & 0 &  45 & 80.6  & 1 \\
CA125 & 1 &  50 &  17.5 & 0 &  41 & 80.4  & 1 \\
LH    & 2 &  50 &  17.5 & 0 &  41 & 81.9  & 1 \\
PTH   & 1 &  50 &  17.1 & 0 &  74 & 102.1 & 1 \\
COR   & 1 &  50 &  17.0 & 0 &  35 & 79.1  & 1 \\
T3L   & 1 &  50 &  15.7 & 0 &  37 & 85.6  & 1 \\
PRL   & 1 &  50 &  13.6 & 0 &  32 & 81.3  & 1 \\
CA199 & 2 &  50 &  13.6 & 0 &  28 & 79.3  & 1 \\
TPO   & 2 &  50 &  12.8 & 0 &  36 & 85.4  & 1 \\
PSAL  & 1 &  50 &  12.6 & 0 &  32 & 85.3  & 1 \\
EST   & 2 &  50 &  12.3 & 0 &  28 & 85.0  & 1 \\
PROG  & 2 &  50 &   8.8 & 0 &  22 & 86.5  & 1 \\
INS   & 1 &  50 &   6.7 & 0 &  17 & 88.4  & 1 \\
AFP   & 1 &  50 &   6.2 & 0 &  15 & 86.3  & 1 \\
HCG   & 1 &  50 &   2.7 & 0 &   8 & 100.1 & 1 \\
TIRO  & 1 &  50 &   2.5 & 0 &  10 & 102.6 & 1 \\
PEPC  & 1 &  50 &   2.5 & 0 &   8 & 103.1 & 1 \\
ATG   & 2 &  50 &   1.5 & 0 &   6 & 111.2 & 1 \\
SHBG  & 2 &  50 &   1.2 & 0 &   4 & 122.0 & 1 \\
\bottomrule
\multicolumn{8}{l}{$^*$Petacas mínimas en un único analizador: $\lceil d_i^{\max}/c_i \rceil$.
Total mínimo: 57 de 144.}
\end{tabular}
\end{table}

El modelo utiliza el \textbf{máximo diario observado} como escenario de planificación
conservador: la distribución resultante puede absorber el peor día registrado sin recargar
reactivos sobre la marcha. El coeficiente de variación elevado (78--122\,\%) refleja la
inclusión de festivos con volumen muy reducido; usar la media infraestimaría la carga
de un día laboral normal.

\subsection{Equilibrio de carga por tipo de paso}

\begin{table}[H]
\centering
\caption{Carga total máxima por tipo de paso.}
\begin{tabular}{lrrr}
\toprule
\textbf{Tipo} & \textbf{Pruebas} & \textbf{Demanda total máx.} & \textbf{Tiempo mínimo$^*$} \\
\midrule
1 paso (450\,det/h) & 15 & $\approx$1\,623\,det/día & $\approx$3.61\,h \\
2 pasos (225\,det/h) & 13 & $\approx$1\,096\,det/día & $\approx$4.87\,h \\
\bottomrule
\multicolumn{4}{l}{$^*$Si todas las pruebas del tipo estuviesen en un único analizador.}
\end{tabular}
\end{table}

Con 3 analizadores compartiendo carga: tiempo esperado $\approx \max(1623/1350,\,1096/675) \approx 1.62$\,h.
El makespan óptimo obtenido es \textbf{172.2\,min (2.87\,h)}, dentro de la jornada de 8\,h
(35.9\,\% de utilización).

\subsection{Co-ocurrencias}

Las co-ocurrencias miden cuántas peticiones mensuales se verían penalizadas si un par de
pruebas está repartido entre distintos equipos.

\paragraph{El efecto ``Big Six''.}
Los 10 pares más frecuentes del mes están completamente dominados por las 6 pruebas de
mayor volumen --- TSH, VITD, B12, FOL, T4L y PSA (\textbf{Big Six}):

\begin{table}[H]
\centering
\caption{Top 10 pares co-ocurrentes globales (enero 2026).}
\begin{tabular}{lr}
\toprule
\textbf{Par} & \textbf{Co-oc./mes} \\
\midrule
TSH + VITD & 3\,954 \\
FOL + B12  & 3\,758 \\
TSH + B12  & 3\,659 \\
TSH + T4L  & 3\,609 \\
FOL + TSH  & 2\,901 \\
B12 + VITD & 2\,458 \\
FOL + VITD & 1\,782 \\
PSA + TSH  & 1\,652 \\
T4L + B12  & 1\,057 \\
T4L + VITD &   989  \\
\bottomrule
\end{tabular}
\end{table}

Esta dominancia enmascara la estructura de co-ocurrencia del resto de pruebas. Por ello se
aplica un tratamiento diferenciado:
\begin{itemize}
  \item \textbf{Big Six $\to$ restricción de presencia (R6b):} se fuerza $x_{ij}=1$ para
        toda $i \in I_B$ y todo $j \in J$. Las 6 pruebas están siempre en los 3
        analizadores, garantizando automáticamente la co-ubicación de cualquier par del
        Big Six en todos los equipos.
  \item \textbf{Resto de pruebas $\to$ objetivo de cohesión (Fase 2):} los grupos
        co-ocurrentes entre pruebas no Big Six se incluyen como parámetros $f_g^{(\ell)}$
        y se maximizan en la función objetivo multi-nivel.
\end{itemize}

Eliminando el Big Six del análisis, emerge la estructura clínica real: un
\textbf{grupo hormonal reproductor} (FSH, LH, EST, PRL, TES, PROG) y un
\textbf{grupo oncológico} (CEA, CA125, CA153, CA199) con fuertes co-solicitudes internas.

\begin{table}[H]
\centering
\caption{Top 10 pares co-ocurrentes excluyendo Big Six (enero 2026).}
\begin{tabular}{lrr}
\toprule
\textbf{Par} & \textbf{Co-oc./mes} & \textbf{Grupo clínico} \\
\midrule
CEA + CA153  & 590 & Oncológico \\
FSH + LH     & 500 & Hormonal   \\
CEA + CA199  & 381 & Oncológico \\
FSH + EST    & 368 & Hormonal   \\
LH + EST     & 336 & Hormonal   \\
LH + PRL     & 305 & Hormonal   \\
TES + LH     & 305 & Hormonal   \\
FSH + PRL    & 301 & Hormonal   \\
FSH + TES    & 277 & Hormonal   \\
FSH + PROG   & 258 & Hormonal   \\
\bottomrule
\end{tabular}
\end{table}

El modelo trabaja con \textbf{50 grupos} en total: los \textbf{10 más frecuentes por cada
nivel} de agrupación (parejas, tripletas, cuartetos, quintuplos y sextuplos). El quíntuplo
más frecuente del panel hormonal es:
\[
  \{\text{EST},\,\text{FSH},\,\text{LH},\,\text{PROG},\,\text{TES}\}
  \quad\longrightarrow\quad 139\text{ co-ocurrencias/mes.}
\]

\clearpage
% ============================================================
\section{Formulación matemática del modelo}
% ============================================================

\subsection{Conjuntos e índices}

\begin{table}[H]
\centering
\caption{Conjuntos del modelo.}
\begin{tabularx}{\linewidth}{lX}
\toprule
\textbf{Símbolo} & \textbf{Descripción} \\
\midrule
$I$ & Conjunto de 28 pruebas clínicas. \\
$J$ & Conjunto de 3 analizadores: $\{\text{DxI1},\text{DxI2},\text{DxI3}\}$. \\
$I_E \subset I$ & Pruebas exentas de contingencia (bajo volumen): $\{\text{HCG, TIRO, PEPC, ATG, SHBG}\}$. \\
$I_1 \subset I$ & Pruebas de 1 paso: $\{\text{TSH, VITD, PSA, PSAL, CA125, AFP, PTH, T3L, PRL, HCG, TES, INS, PEPC, COR, TIRO}\}$. \\
$I_B \subset I$ & Big Six (forzadas en los 3 analizadores): $\{\text{TSH, VITD, B12, FOL, T4L, PSA}\}$. \\
$\mathcal{G}^{(\ell)}$ & Grupos de nivel $\ell$, con $\ell \in \{2,3,4,5,6\}$. $|\mathcal{G}^{(\ell)}| = 10$ (top 10 por frecuencia). \\
$C_g^{(\ell)} \subset I \setminus I_B$ & Pruebas que componen el grupo $g \in \mathcal{G}^{(\ell)}$. $|C_g^{(\ell)}| = \ell$. \\
\bottomrule
\end{tabularx}
\end{table}

\subsection{Parámetros}

\begin{table}[H]
\centering
\caption{Parámetros del modelo.}
\begin{tabularx}{\linewidth}{llX}
\toprule
\textbf{Símbolo} & \textbf{Unidades} & \textbf{Descripción} \\
\midrule
$\text{cap}_j$ & posiciones & Posiciones de reactivo disponibles en el analizador $j$. \\
$c_i$ & det/petaca & Determinaciones por petaca de la prueba $i$. \\
$v_i$ & det/h & Velocidad máxima de procesamiento. \\
$d_i$ & det/día & Demanda máxima diaria. \\
$H$ & h & Horas de trabajo disponibles por día. \\
$\delta$ & --- & Diferencia máxima de pruebas de 1 paso entre analizadores. \\
$f_g^{(\ell)}$ & peticiones/mes & Co-ocurrencias mensuales del grupo $g$ de nivel $\ell$. \\
$\tau$ & det/día & Umbral de bajo volumen para R6c ($\tau = 20$). \\
$T^*$ & h & Makespan óptimo de la Fase 1 (fija la cota en Fase 2). \\
\bottomrule
\end{tabularx}
\end{table}

\textbf{Valores de los parámetros de infraestructura:}
\[
  \text{cap}_{\text{DxI1}} = 50, \qquad
  \text{cap}_{\text{DxI2}} = 47, \qquad
  \text{cap}_{\text{DxI3}} = 47
\]
\[
  H = 8\text{\,h/día}, \qquad \delta = 3
\]
\[
  c_{\text{TSH}} = 200, \qquad c_{\text{PSA}} = 100, \qquad
  c_i = 50 \quad \forall\, i \notin \{\text{TSH},\,\text{PSA}\}
\]

\subsection{Variables de decisión}

\begin{table}[H]
\centering
\caption{Variables de decisión del modelo.}
\begin{tabularx}{\linewidth}{llX}
\toprule
\textbf{Variable} & \textbf{Tipo} & \textbf{Descripción} \\
\midrule
$p_{ij} \in \mathbb{Z}_{\geq 0}$ & Entera & Petacas de la prueba $i$ cargadas en el analizador $j$. \\
$x_{ij} \in \{0,1\}$ & Binaria & 1 si la prueba $i$ está disponible en el analizador $j$. \\
$y_{ij} \in [0,1]$ & Continua & Fracción de la demanda de $i$ procesada en $j$. \\
$T_{\max} \geq 0$ & Continua & Makespan: tiempo del analizador más cargado (horas). \\
$w_g^{(\ell)}[j] \in \{0,1\}$ & Binaria & 1 si \textbf{todas} las pruebas del grupo $g \in \mathcal{G}^{(\ell)}$ están en $j$. \\
\bottomrule
\end{tabularx}
\end{table}

\textbf{Relaciones entre variables:}
$x_{ij} = 1 \Rightarrow p_{ij} \geq 1$;\ \
$x_{ij} = 0 \Rightarrow p_{ij} = 0,\, y_{ij} = 0$;\ \
$y_{ij} > 0$ solo si $x_{ij} = 1$.

\subsection{Función objetivo (optimización lexicográfica en dos fases)}

\paragraph{Fase 1 — Minimizar el makespan:}
\[
  \min \quad T_{\max}
\]
Minimizar el tiempo del analizador más cargado equivale a \textbf{equilibrar la carga}
entre los tres equipos.

\paragraph{Fase 2 — Maximizar la cohesión jerárquica:}
\[
  \max \quad \sum_{\ell=2}^{6}\ \sum_{g \in \mathcal{G}^{(\ell)}}\ \sum_{j \in J}
             f_g^{(\ell)} \cdot w_g^{(\ell)}[j]
\]
Para cada grupo $g$ de nivel $\ell$ y analizador $j$, se suma $f_g^{(\ell)}$ cada vez
que todas las pruebas del grupo están juntas en $j$. La restricción R11
($T_{\max} \leq T^*$) garantiza que la búsqueda de cohesión no empeora el makespan.

Usar grupos de hasta 6 pruebas permite capturar directamente el valor de mantener paneles
clínicos completos co-ubicados: si el panel hormonal $\{\text{EST,FSH,LH,PROG,PRL,TES}\}$
queda en el mismo analizador, el beneficio se contabiliza en cinco niveles de agrupación
simultáneamente.

\subsection{Restricciones}

\paragraph{R1 — Capacidad de posiciones:}
\[
  \sum_{i \in I} p_{ij} \leq \text{cap}_j \qquad \forall\, j \in J
\]

\paragraph{R2 — Petacas suficientes para la demanda enrutada:}
\[
  p_{ij} \cdot c_i \geq y_{ij} \cdot d_i \qquad \forall\, i \in I,\ \forall\, j \in J
\]

\paragraph{R3 — Cobertura total de la demanda:}
\[
  \sum_{j \in J} y_{ij} = 1 \qquad \forall\, i \in I
\]

\paragraph{R4 — Enrutamiento solo si la prueba está instalada:}
\[
  y_{ij} \leq x_{ij} \qquad \forall\, i \in I,\ \forall\, j \in J
\]

\paragraph{R5 — Coherencia entre petacas y presencia:}
\[
  p_{ij} \leq \text{cap}_j \cdot x_{ij} \qquad \forall\, i \in I,\ \forall\, j \in J
\]
\[
  x_{ij} \leq p_{ij} \qquad \forall\, i \in I,\ \forall\, j \in J
\]
La primera garantiza que si $x_{ij}=0$ entonces $p_{ij}=0$; la segunda que si
$x_{ij}=1$ entonces $p_{ij} \geq 1$.

\paragraph{R6 — Plan de contingencia:}
\[
  \sum_{j \in J} x_{ij} \geq 2 \qquad \forall\, i \in I \setminus I_E
\]
\[
  \sum_{j \in J} x_{ij} \geq 1 \qquad \forall\, i \in I_E
\]
Las pruebas no exentas deben estar en al menos dos analizadores. Las pruebas de muy bajo
volumen ($I_E$) quedan exentas: duplicarlas no es coste-efectivo.

\paragraph{R6b — Big Six en los tres analizadores:}
\[
  x_{ij} = 1 \qquad \forall\, i \in I_B,\ \forall\, j \in J
\]
Las seis pruebas de mayor volumen están forzadas en los 3 analizadores, garantizando su
co-ubicación universal por restricción y liberando el objetivo de Fase 2 para cohesionar
el resto del panel.

\paragraph{R6c — Pruebas de bajo volumen en máximo dos analizadores:}
\[
  \sum_{j \in J} x_{ij} \leq 2 \qquad \forall\, i \in I \setminus I_B :\ d_i < \tau
\]
donde $\tau = 20$\,det/día. El coste acumulado de calibración en un tercer equipo supera el
beneficio para pruebas de muy bajo volumen. Pruebas afectadas: INS (17), AFP (15),
HCG (8), TIRO (10), PEPC (8), ATG (6), SHBG (4\,det/día).

\paragraph{R7 — Sincronización PSA / PSAL:}
\[
  x_{\text{PSAL},j} \leq x_{\text{PSA},j} \qquad \forall\, j \in J
\]
PSAL solo puede estar en un analizador donde PSA también esté disponible. PSA puede estar
sin PSAL (hay peticiones que no solicitan la fracción libre).

\paragraph{R8 — Definición del makespan:}
\[
  \sum_{i \in I} \frac{y_{ij} \cdot d_i}{v_i} \leq T_{\max} \qquad \forall\, j \in J
\]

\paragraph{R9 — Equilibrio de pruebas de 1 paso:}
\[
  \left|\sum_{i \in I_1} x_{ij} - \sum_{i \in I_1} x_{ik}\right| \leq \delta
  \qquad \forall\, j,k \in J,\ j < k
\]
Evita que un analizador acumule pruebas de 2 pasos mientras otro acumula las de 1 paso.

\paragraph{R10 — Jornada laboral:}
\[
  T_{\max} \leq H
\]

\paragraph{R11 — Fijación del makespan entre fases:}
\[
  T_{\max} \leq T^* \qquad (\text{activa solo en Fase 2})
\]
En Fase 1 esta restricción es inactiva ($T^* = +\infty$). Al inicio de la Fase 2 se asigna
$T^* \leftarrow T_{\max}^{\text{Fase 1}}$.

\paragraph{R15--R16 — Linealización de $w_g^{(\ell)}$ (AND de $n$ binarias):}

La variable $w_g^{(\ell)}[j]$ debe ser 1 si y solo si todas las $\ell$ pruebas del grupo
$g$ están instaladas en $j$:
\[
  w_g^{(\ell)}[j] = \bigwedge_{i \in C_g^{(\ell)}} x_{ij}
\]

El AND lógico se linealiza con dos familias estándar:

\textbf{R15 — Cota superior} (si falta alguna prueba del grupo, $w = 0$):
\[
  w_g^{(\ell)}[j] \leq x_{ij}
  \qquad \forall\, g \in \mathcal{G}^{(\ell)},\ \forall\, i \in C_g^{(\ell)},\
  \forall\, j \in J,\ \forall\, \ell
\]

\textbf{R16 — Cota inferior} (si todas las pruebas están, $w$ debe ser 1):
\[
  w_g^{(\ell)}[j] \geq \sum_{i \in C_g^{(\ell)}} x_{ij} - (\ell - 1)
  \qquad \forall\, g \in \mathcal{G}^{(\ell)},\ \forall\, j \in J,\ \forall\, \ell
\]

Esta formulación generaliza la relajación de McCormick (válida solo para $\ell = 2$) a
grupos de cualquier tamaño con solo dos familias de restricciones, frente a las
$2^{\ell}-1$ inecuaciones de la envoltura convexa completa.

\subsection{Formulación compacta}

\[
  \underset{\text{Fase 1}}{\min}\ T_{\max}
  \quad \longrightarrow \quad
  \underset{\text{Fase 2}}{\max}
  \sum_{\ell=2}^{6}\sum_{g \in \mathcal{G}^{(\ell)}}\sum_{j\in J} f_g^{(\ell)}\,w_g^{(\ell)}[j]
\]

\[
\text{s.a.} \quad
\begin{dcases}
\displaystyle\sum_{i} p_{ij} \leq \text{cap}_j & \forall j \\[6pt]
p_{ij}\, c_i \geq y_{ij}\, d_i & \forall i,j \\[6pt]
\displaystyle\sum_{j} y_{ij} = 1 & \forall i \\[6pt]
y_{ij} \leq x_{ij} & \forall i,j \\[6pt]
p_{ij} \leq \text{cap}_j\, x_{ij},\quad x_{ij} \leq p_{ij} & \forall i,j \\[6pt]
\displaystyle\sum_{j} x_{ij} \geq 2 & \forall i \notin I_E \\[6pt]
\displaystyle\sum_{j} x_{ij} \geq 1 & \forall i \in I_E \\[6pt]
x_{ij} = 1 & \forall i \in I_B,\ \forall j \\[6pt]
\displaystyle\sum_{j} x_{ij} \leq 2 & \forall i \notin I_B :\ d_i < \tau \\[6pt]
x_{\text{PSAL},j} \leq x_{\text{PSA},j} & \forall j \\[6pt]
\displaystyle\sum_{i} \dfrac{y_{ij} d_i}{v_i} \leq T_{\max} & \forall j \\[6pt]
\left|\displaystyle\sum_{i \in I_1} x_{ij} - \sum_{i \in I_1} x_{ik}\right| \leq \delta & \forall j < k \\[6pt]
T_{\max} \leq H \\[6pt]
T_{\max} \leq T^* \\[6pt]
w_g^{(\ell)}[j] \leq x_{ij} & \forall \ell,\ g,\ i \in C_g^{(\ell)},\ j \\[6pt]
w_g^{(\ell)}[j] \geq \displaystyle\sum_{i \in C_g^{(\ell)}} x_{ij} - (\ell-1) & \forall \ell,\ g,\ j \\[6pt]
p_{ij} \in \mathbb{Z}_{\geq 0},\quad x_{ij},\,w_g^{(\ell)}[j] \in \{0,1\},\quad y_{ij} \in [0,1],\quad T_{\max} \geq 0
\end{dcases}
\]

\subsection{Clasificación del modelo}

El modelo es un \textbf{MILP} (\textit{Mixed-Integer Linear Program}): variables continuas
($y_{ij}$, $T_{\max}$), enteras ($p_{ij}$) y binarias ($x_{ij}$, $w_g^{(\ell)}[j]$),
con todas las restricciones y objetivos \textbf{lineales}.

\begin{table}[H]
\centering
\caption{Dimensión del problema ($|I|=28$, $|J|=3$, $|\mathcal{G}^{(\ell)}|=10$).}
\begin{tabular}{lrr}
\toprule
\textbf{Elemento} & \textbf{Cálculo} & \textbf{Cantidad} \\
\midrule
Var. binarias $x_{ij}$ & $28 \times 3$ & 84 \\
Var. binarias $w_g^{(\ell)}[j]$ & $5 \times 10 \times 3$ & 150 \\
Var. enteras $p_{ij}$ & $28 \times 3$ & 84 \\
Var. continuas $y_{ij}$ y $T_{\max}$ & $28 \times 3 + 1$ & 85 \\
\textbf{Total variables} & & \textbf{403} \\
\midrule
R1 & $3$ & 3 \\
R2 & $28 \times 3$ & 84 \\
R3 & $28$ & 28 \\
R4 & $28 \times 3$ & 84 \\
R5 (R5a + R5b) & $2 \times 28 \times 3$ & 168 \\
R6 & $28$ & 28 \\
R6b & $6 \times 3$ & 18 \\
R6c & 7 pruebas con $d_i < 20$, excluido $I_B$ & 7 \\
R7 & $3$ & 3 \\
R8 & $3$ & 3 \\
R9 & $2 \times \binom{3}{2}$ & 6 \\
R10 & 1 & 1 \\
R11 & 1 & 1 \\
R15 — cota superior $w$ & $10 \times (2{+}3{+}4{+}5{+}6) \times 3$ & 600 \\
R16 — cota inferior $w$ & $5 \times 10 \times 3$ & 150 \\
\textbf{Total restricciones} & & \textbf{1\,188} \\
\bottomrule
\end{tabular}
\end{table}

Reducir a top 10 por nivel recorta el modelo casi a la mitad en variables de co-ubicación
(300 $\to$ 150) y restricciones de linealización (1\,500 $\to$ 750), manteniendo la
representatividad al conservar únicamente los grupos con mayor impacto clínico.

\subsection{Nota sobre la variable de enrutamiento $y_{ij}$}

La variable $y_{ij}$ es una \textbf{abstracción de planificación}: representa la fracción
de la demanda diaria de la prueba $i$ que se diseña para el analizador $j$. En la
operativa real, esta distribución emerge del enrutamiento automático de la cadena. El modelo
usa $y_{ij}$ para calcular la carga de forma lineal, evitando la no-linealidad que
aparecería si se modelase el reparto proporcional a petacas.

\clearpage
% ============================================================
\section{Instrucciones de uso (AMPL)}
% ============================================================

\subsection{Archivos del modelo}

\begin{table}[H]
\centering
\caption{Archivos del proyecto AMPL.}
\begin{tabularx}{\linewidth}{lX}
\toprule
\textbf{Archivo} & \textbf{Contenido} \\
\midrule
\texttt{modelo\_dxi.mod} & Definición del modelo AMPL: conjuntos, parámetros, variables, objetivos y restricciones. \\
\texttt{datos\_dxi.dat} & Parámetros numéricos: infraestructura, demandas y los 50 grupos de co-ocurrencia (top 10 por nivel, niveles 2--6). \\
\texttt{run\_dxi.run} & Script de ejecución: resuelve las dos fases, imprime resultados y ejecuta el postproceso proporcional. El flujo completo es autocontenido. \\
\bottomrule
\end{tabularx}
\end{table}

\texttt{postproceso.py} es un script Python auxiliar que implementa la misma lógica de
distribución proporcional como referencia para verificación independiente, pero
\textbf{no es necesario ejecutarlo}: \texttt{run\_dxi.run} ya incorpora el postproceso.

\subsection{Ejecución}

\begin{lstlisting}[language={},caption={Desde la consola interactiva de AMPL.}]
ampl: include run_dxi.run;
\end{lstlisting}

\begin{lstlisting}[language=bash,caption={Desde terminal (línea de comandos).}]
ampl run_dxi.run
\end{lstlisting}

\subsection{Cambiar el solver}

En \texttt{run\_dxi.run}, modificar la línea:
\begin{lstlisting}[language={}]
option solver cplex;   # -> glpk, highs, gurobi, cbc...
\end{lstlisting}

\subsection{Demanda usada en el modelo}

El archivo \texttt{datos\_dxi.dat} utiliza los \textbf{máximos diarios observados} como
valores del parámetro \texttt{dem}. Para experimentar con la demanda media, sustituir el
bloque \texttt{param dem} por las medias de la tabla~4.1.

\subsection{Modificar el parámetro de equilibrio $\delta$}

\begin{lstlisting}[language={}]
param delta := 2;   # mayor equilibrio de pasos 1/2
param delta := 5;   # mayor flexibilidad
\end{lstlisting}

\subsection{Postproceso: llenado proporcional de posiciones}

Tras las dos fases de AMPL, las posiciones sobrantes se distribuyen mediante el siguiente
algoritmo:

\textbf{Paso 1 — Mínimo operativo garantizado.}
Para cada prueba $i$ instalada en el analizador $j$:
\[
  p_{ij}^{\min} = \max\!\left(1,\ \left\lceil \frac{y_{ij} \cdot d_i}{c_i} \right\rceil\right)
\]

\textbf{Paso 2 — Reparto proporcional de huecos sobrantes.}
Sea $E_j = \text{cap}_j - \sum_i p_{ij}^{\min}$ las posiciones libres tras el paso 1.
Cada prueba instalada con $y_{ij} > 0$ recibe:
\[
  \text{extra}_{ij} = E_j \cdot
  \frac{y_{ij} \cdot d_i}{\displaystyle\sum_{\substack{k:\,x_{kj}=1 \\ y_{kj}>0}} y_{kj} \cdot d_k}
\]
Los valores se redondean por la parte entera inferior; los residuos se asignan a las pruebas
con mayor parte fraccionaria. En caso de empate, el desempate se resuelve por orden de
aparición en el conjunto \texttt{PRUEBAS} (\texttt{ord()}), garantizando un orden total y
evitando sobreasignación. El resultado satisface exactamente $\sum_i p_{ij} = \text{cap}_j$.

\textbf{Justificación:} las pruebas con mayor demanda reciben más petacas extra, de modo que
si un analizador cae, los restantes tienen suficiente capacidad de reactivo para absorber
el trabajo sin recargas adicionales.

\subsection{Interpretación de los resultados}

El flujo completo (AMPL + postproceso) produce los siguientes bloques de salida:

\textbf{Fase 1 — Makespan mínimo:}
\begin{enumerate}
  \item \textbf{Estado de la solución} y \textbf{makespan óptimo} en horas y minutos.
  \item \textbf{Carga por analizador:} posiciones usadas, tiempo de procesamiento y
        porcentaje de la jornada.
  \item \textbf{Asignación detallada:} para cada analizador, qué pruebas están instaladas,
        cuántas petacas tienen, qué fracción de demanda gestionan y cuánto tiempo consumen.
  \item \textbf{Tabla de contingencia:} en cuántos analizadores está cada prueba (columna
        \texttt{nDxI}) y petacas en cada uno.
  \item \textbf{Balance de pasos:} distribución de pruebas de 1 y 2 pasos entre analizadores.
  \item \textbf{Uso total de posiciones:} posiciones usadas vs.\ disponibles vs.\ libres.
  \item \textbf{Verificación PSA/PSAL:} comprueba que PSAL no está en ningún analizador
        donde PSA esté ausente (condición $x_{\text{PSAL},j} \leq x_{\text{PSA},j}$).
\end{enumerate}

\textbf{Fase 2 — Cohesión máxima:}
\begin{enumerate}[resume]
  \item \textbf{Puntuación de cohesión:} co-ocurrencias mensuales cubiertas por un único
        analizador, expresadas como valor absoluto y porcentaje sobre el máximo teórico.
  \item \textbf{Tabla de pares co-ubicados:} para cada par $(i,k)$, indica en qué
        analizador(es) quedaron juntos o si resultaron fragmentados.
  \item \textbf{Cohesión por analizador:} número de pares co-ubicados y score parcial
        de cada DxI.
\end{enumerate}

\textbf{Postproceso proporcional:}
\begin{enumerate}[resume]
  \item \textbf{Distribución final de petacas:} tabla por analizador con petacas mínimas,
        petacas finales, posiciones extra añadidas y buffer de capacidad respecto a la
        demanda máxima.
\end{enumerate}

\clearpage
% ============================================================
\section{Resultados de la solución óptima}
% ============================================================

\subsection{Comparación con la distribución actual del laboratorio}

\begin{table}[H]
\centering
\caption{Comparación entre la distribución empírica actual y la solución óptima.}
\begin{tabularx}{\linewidth}{Xcc>{\raggedright\arraybackslash}X}
\toprule
\textbf{Métrica} & \textbf{Actual} & \textbf{Óptima} & \textbf{Mejora} \\
\midrule
Makespan (dem. media)$^\dagger$ & $\approx$91.8\,min & \textbf{83.8\,min} & $-8$\,min ($-8.7$\,\%) \\
Balance de carga (media) & 91.8/90.5/69.2\,min & 83.8/83.8/83.8\,min & Desv. $22.6 \to \mathbf{0}$\,min \\
Makespan (dem. máxima) & no optimizada & \textbf{172.2\,min} & --- \\
R6b Big Six en 3 máquinas & Cumplida & Garantizada por R6b & Formalizada \\
R7 PSA/PSAL & Compatible & Cumplida & --- \\
R6 Contingencia INS & \textbf{VIOLACIÓN} (1 máquina) & Cumplida (2 máquinas) & Corregida \\
Cohesión de peticiones$^\ddagger$ & Sin garantía formal & \textbf{99.7\,\%} (25\,804/25\,872) & Optimizada \\
\bottomrule
\end{tabularx}
\end{table}

\begin{small}
$^\dagger$\,\textit{Comparación de makespan:} las filas de demanda media comparan ambas
distribuciones en igualdad de condiciones (mismo escenario de carga). El modelo de
producción usa demanda máxima (172.2\,min), pero la distribución actual no fue diseñada
bajo ese escenario. La mejora de balance (desvío 22.6 $\to$ 0\,min) es válida en ambos
escenarios.

$^\ddagger$\,\textit{Cohesión:} el score 25\,804/25\,872 = 99.7\,\% integra los top 10
grupos de cinco niveles. Parejas, tripletas y cuartetos alcanzan el 100\,\%. La pérdida
del 0.3\,\% se concentra en tres sextuplos (SEXT\_8/9/10, que contienen INS) y un
quintuplo (QUIN\_7, que contiene AFP), limitados por R6c a solo 2 analizadores.
\end{small}

\subsection{Resultados del modelo (demanda máxima)}

\textbf{Solver:} HiGHS 1.11.0 --- solución óptima en 217 iteraciones símplex (Fase 1)
$+$ 195 (Fase 2), 1 nodo de ramificación en cada fase.

\medskip
\textbf{Fase 1 — Makespan óptimo:} $\mathbf{2.8696\text{ h} = 172.2\text{ min}}$
(35.9\,\% de la jornada de 8\,h). Los tres analizadores quedan exactamente igualados ---
\textbf{balance perfecto}.

\begin{table}[H]
\centering
\caption{Carga por analizador (Fase 1).}
\begin{tabular}{cccc}
\toprule
\textbf{Analizador} & \textbf{Pos. usadas} & \textbf{Carga} & \textbf{\% Jornada} \\
\midrule
DxI1 & 23 & 172.2\,min & 35.9\,\% \\
DxI2 & 27 & 172.2\,min & 35.9\,\% \\
DxI3 & 25 & 172.2\,min & 35.9\,\% \\
\bottomrule
\end{tabular}
\end{table}

\textbf{Fase 2 — Cohesión: 25\,804 / 25\,872 = 99.7\,\%.}
Parejas (10\,863), tripletas (6\,990) y cuartetos (4\,896) alcanzan el 100\,\% --- todos
sus grupos quedan co-ubicados en los \textbf{tres analizadores}. La pérdida del 0.3\,\%
afecta solo a:
\begin{itemize}
  \item Tres sextuplos (SEXT\_8, SEXT\_9, SEXT\_10), que contienen INS (limitado a 2
        analizadores por R6c).
  \item El quintuplo QUIN\_7 (AFP CA125 CA153 CA199 CEA), cuyo componente AFP está
        limitado a 2 analizadores por R6c.
\end{itemize}

\textbf{Verificación PSA/PSAL (Fase 2):}
DxI1: PSA=1 PSAL=1 $\to$ OK $\mid$ DxI2: PSA=1 PSAL=0 $\to$ OK $\mid$
DxI3: PSA=1 PSAL=1 $\to$ OK.
La restricción $x_{\text{PSAL},j} \leq x_{\text{PSA},j}$ se cumple en todos los
analizadores.

\subsection{Distribución de pruebas por número de analizadores}

El resultado final (Fase 2) instala 16 pruebas en los 3 analizadores, 7 en 2 y 5 en 1:

\begin{table}[H]
\centering
\caption{Asignación final de pruebas por analizador (Fase 2).}
\small
\begin{tabular}{lrccccc}
\toprule
\textbf{Prueba} & \textbf{Dem. máx} & \textbf{DxI1} & \textbf{DxI2} & \textbf{DxI3} & \textbf{N.º máq.} & \textbf{Nota} \\
\midrule
TSH   & 765 & \cmark & \cmark & \cmark & 3 & Big Six \\
VITD  & 322 & \cmark & \cmark & \cmark & 3 & Big Six \\
B12   & 300 & \cmark & \cmark & \cmark & 3 & Big Six \\
FOL   & 260 & \cmark & \cmark & \cmark & 3 & Big Six \\
T4L   & 240 & \cmark & \cmark & \cmark & 3 & Big Six \\
PSA   & 149 & \cmark & \cmark & \cmark & 3 & Big Six \\
CEA   &  88 & \cmark & \cmark & \cmark & 3 & Cohesión \\
FSH   &  45 & \cmark & \cmark & \cmark & 3 & Cohesión \\
TES   &  43 & \cmark & \cmark & \cmark & 3 & Cohesión \\
CA153 &  45 & \cmark & \cmark & \cmark & 3 & Cohesión \\
LH    &  41 & \cmark & \cmark & \cmark & 3 & Cohesión \\
COR   &  35 & \cmark & \cmark & \cmark & 3 & Cohesión \\
PRL   &  32 & \cmark & \cmark & \cmark & 3 & Cohesión \\
CA199 &  28 & \cmark & \cmark & \cmark & 3 & Cohesión \\
EST   &  28 & \cmark & \cmark & \cmark & 3 & Cohesión \\
PROG  &  22 & \cmark & \cmark & \cmark & 3 & Cohesión \\
\midrule
CA125 &  41 & ---    & \cmark & \cmark & 2 & \\
PTH   &  74 & \cmark & \cmark & ---    & 2 & \\
T3L   &  37 & \cmark & ---    & \cmark & 2 & \\
TPO   &  36 & \cmark & \cmark & ---    & 2 & \\
PSAL  &  32 & \cmark & ---    & \cmark & 2 & R7: PSAL $\subseteq$ PSA \\
INS   &  17 & \cmark & ---    & \cmark & 2 & R6c: dem $<$ 20 \\
AFP   &  15 & ---    & \cmark & \cmark & 2 & R6c: dem $<$ 20 \\
\midrule
HCG   &   8 & \cmark & ---    & ---    & 1 & R6c, exenta \\
TIRO  &  10 & ---    & ---    & \cmark & 1 & R6c, exenta \\
PEPC  &   8 & ---    & \cmark & ---    & 1 & R6c, exenta \\
ATG   &   6 & ---    & ---    & \cmark & 1 & R6c, exenta \\
SHBG  &   4 & ---    & ---    & \cmark & 1 & R6c, exenta \\
\bottomrule
\end{tabular}
\end{table}

\textbf{Patrón observado:} las 6 pruebas del Big Six están en los 3 analizadores por R6b.
La Fase 2 expande a 3 analizadores 10 pruebas adicionales: el panel hormonal reproductor
completo (FSH, LH, TES, PRL, COR, EST, PROG) más CEA, CA153 y CA199 del grupo oncológico.
CA125 solo aparece en QUIN\_7 (AFP CA125 CA153 CA199 CEA, $f=37$), cuyo componente AFP
está limitado a 2 analizadores por R6c; en consecuencia CA125 queda en solo 2 máquinas.
Las 7 pruebas de muy bajo volumen se distribuyen entre 1 y 2 máquinas: INS y AFP en 2;
HCG, TIRO, PEPC, ATG y SHBG en 1.

\subsection{Distribución final de petacas (postproceso proporcional, 144 posiciones)}

Tras el postproceso, las 144 posiciones quedan completamente ocupadas. La columna
\emph{Factor} indica $\text{Cap. total}/d_i$; un factor $\geq 2$ garantiza que los
analizadores restantes pueden absorber la carga si uno falla.

\begin{table}[H]
\centering
\caption{Distribución final de petacas por prueba y analizador.}
\small
\begin{tabular}{lrcccrc}
\toprule
\textbf{Prueba} & \textbf{Dem. máx} & \textbf{DxI1} & \textbf{DxI2} & \textbf{DxI3} & \textbf{Cap. total} & \textbf{Factor} \\
\midrule
TSH   & 765 &  7 &  6 &  8 & 4\,200 & 5.5$\times$ \\
VITD  & 322 &  2 &  2 &  9 &   650 & 2.0$\times$ \\
B12   & 300 & 10 &  2 &  2 &   700 & 2.3$\times$ \\
FOL   & 260 &  5 &  7 &  1 &   650 & 2.5$\times$ \\
T4L   & 240 &  2 &  7 &  2 &   550 & 2.3$\times$ \\
PSA   & 149 &  4 &  2 &  1 &   700 & 4.7$\times$ \\
CEA   &  88 &  1 &  2 &  2 &   250 & 2.8$\times$ \\
FSH   &  45 &  2 &  1 &  1 &   200 & 4.4$\times$ \\
TES   &  43 &  2 &  1 &  1 &   200 & 4.7$\times$ \\
CA153 &  45 &  1 &  1 &  2 &   200 & 4.4$\times$ \\
CA125 &  41 & --- &  1 &  2 &   150 & 3.7$\times$ \\
LH    &  41 &  1 &  1 &  2 &   200 & 4.9$\times$ \\
PTH   &  74 &  2 &  2 & --- &  200 & 2.7$\times$ \\
COR   &  35 &  1 &  2 &  1 &   200 & 5.7$\times$ \\
T3L   &  37 &  1 & --- &  2 &  150 & 4.1$\times$ \\
PRL   &  32 &  1 &  2 &  1 &   200 & 6.3$\times$ \\
CA199 &  28 &  1 &  2 &  1 &   200 & 7.1$\times$ \\
TPO   &  36 &  2 &  1 & --- &  150 & 4.2$\times$ \\
PSAL  &  32 &  1 & --- &  2 &  150 & 4.7$\times$ \\
EST   &  28 &  1 &  2 &  1 &   200 & 7.1$\times$ \\
PROG  &  22 &  1 &  1 &  1 &   150 & 6.8$\times$ \\
INS   &  17 &  1 & --- &  1 &  100 & 5.9$\times$ \\
AFP   &  15 & --- &  1 &  1 &   100 & 6.7$\times$ \\
HCG   &   8 &  1 & --- & --- &  50 & 6.3$\times$ \\
TIRO  &  10 & --- & --- &  1 &   50 & 5.0$\times$ \\
PEPC  &   8 & --- &  1 & --- &   50 & 6.3$\times$ \\
ATG   &   6 & --- & --- &  1 &   50 & 8.3$\times$ \\
SHBG  &   4 & --- & --- &  1 &   50 & 12.5$\times$ \\
\midrule
\textbf{TOTAL} & & \textbf{50} & \textbf{47} & \textbf{47} & & \\
\bottomrule
\end{tabular}
\end{table}

Todas las pruebas superan el umbral de contingencia $2\times$; el punto más ajustado es
VITD con exactamente $2.0\times$. T4L alcanza $2.3\times$ y B12 $2.3\times$, ambas
mejoradas respecto a versiones anteriores gracias a las posiciones liberadas por R6c. La
distribución proporcional integrada en el \texttt{.run} garantiza exactamente 144
posiciones ocupadas: DxI1 = 50, DxI2 = 47, DxI3 = 47.

\clearpage
% ============================================================
\section{Conclusiones}
% ============================================================

Este trabajo ha presentado el diseño, formulación y resolución de un modelo de
optimización para la distribución de reactivos en los tres analizadores DxI-9000 de un
laboratorio de análisis clínico. Las principales conclusiones se articulan en torno a
cuatro ejes:

\paragraph{1. Mejora cuantificable sobre la distribución actual.}
La solución óptima reduce el makespan en demanda media en un 8.7\,\% respecto a la
distribución empírica del laboratorio (83.8 frente a 91.8\,min), elimina completamente
el desequilibrio de carga entre analizadores (desvío 22.6 $\to$ 0\,min) y corrige la
única violación de contingencia detectada: INS, que operaba con un solo analizador,
queda ahora garantizada en dos. Adicionalmente, la cohesión de peticiones --- sin
garantía formal en la distribución actual --- alcanza el 99.7\,\% (25\,804/25\,872
co-ocurrencias mensuales cubiertas en un único equipo).

\paragraph{2. Eficacia de la formulación lexicográfica en dos fases.}
La separación jerárquica de objetivos demuestra ser la estrategia correcta para este
problema: la Fase 1 obtiene el makespan mínimo (172.2\,min, 35.9\,\% de la jornada)
con balance perfecto entre los tres analizadores, y la Fase 2 maximiza la cohesión
sin comprometer en absoluto ese resultado. El solver HiGHS 1.11.0 resuelve ambas
fases en un número muy reducido de nodos (1 en cada fase), lo que confirma que la
formulación es ajustada y las restricciones de simetría y equilibrio (R6b, R9)
orientan eficazmente la búsqueda.

\paragraph{3. Valor del tratamiento diferenciado del Big Six y R6c.}
Separar las seis pruebas de mayor volumen (Big Six) mediante una restricción dura en
lugar de incentivarlas en el objetivo libera la Fase 2 para cohesionar el panel
hormonal reproductor y el grupo oncológico, donde reside la variabilidad real de
co-solicitudes. La restricción R6c, que limita a dos analizadores las pruebas con
demanda inferior a 20\,det/día, libera posiciones que se redistribuyen a pruebas de
mayor volumen (T4L pasa a factor 2.3$\times$, B12 a 2.3$\times$, VITD a 2.0$\times$),
reforzando la autonomía del sistema ante un fallo de equipo. El coste de esta
restricción en cohesión es mínimo: únicamente QUIN\_7 y los sextuplos SEXT\_8/9/10,
que contienen AFP e INS respectivamente, pierden la co-ubicación en el tercer
analizador.

\paragraph{4. Linealización AND jerárquica como contribución técnica.}
La generalización de la relajación de McCormick a grupos de $\ell$ variables binarias
mediante solo dos familias de restricciones (R15 y R16) permite capturar directamente
el valor de mantener paneles clínicos completos --- hasta sextuplos --- en el mismo
analizador. Esta formulación escala linealmente con el número de grupos y niveles
($O(\ell \cdot |\mathcal{G}^{(\ell)}| \cdot |J|)$ restricciones), frente a la
complejidad exponencial ($2^{\ell}-1$) de la envoltura convexa completa.

\paragraph{Limitaciones y trabajo futuro.}
El modelo usa la demanda máxima diaria observada como escenario de planificación
conservador, lo que puede sobredimensionar la capacidad en días normales. Una extensión
natural sería incorporar la distribución estocástica de la demanda (p.\,ej., mediante un
modelo robusto o estocástico por escenarios) para obtener soluciones que balanceen
coste de reactivo y riesgo de agotamiento. Asimismo, el modelo asume que los tres
analizadores están siempre operativos; modelar explícitamente los tiempos de
mantenimiento preventivo y las ventanas de recalibración podría refinar las garantías
de contingencia más allá del umbral $2\times$ actualmente asegurado.

% ============================================================
\end{document}
% ============================================================
